\section{Reoptimización y proyección por sitios}
En la parte local se incluyen un factor de capacidad uno por fibra RC marcada y
los factores diagonales activados en $Y$. Los unitarios se pueden mantener; su
corrección de entropía y su determinante normalizado son triviales.
Con $d_i$ incidencias locales por sitio, $q_f=1-\sum_{i\in f}x_i$,
$v_i=x_i(1-x_i)$ y $V_0=\diag(v_i)$, la entropía local de Bethe es
\[
 H_B(x)=\sum_f H_{\rm cat}(x_f)-\sum_i(d_i-1)H_{\rm Ber}(x_i).
\]
Los dos factores RC de cada sitio se incluyen aun si unitarios.
A las actividades $\zeta_i=\exp r_i$ se les da primero libertad en un polidisco
pequeño de radio $\rho_v>0$. En el modelo físico $\zeta_i=O_v(s)$: en el caso
sin diagonales la silla es $x_i=p_{T,i}\le C/j$. La ecuación de estación es
\[
 x_i=\zeta_i \exp(\partial_i R_c(x))
             \frac{\prod_{f\ni i}q_f}{(1-x_i)^{d_i-1}} .
\]
En el punto de torres todas las vacancias marcadas de esta ecuación son $1+O_v(s)$.
La misma ecuación da en el polidisco la silla holomorfa, divisible en cada coordenada
por $\zeta_i$. El jacobiano de la ecuación implícita es identidad en el origen.
En la rama pequeña los términos físicos son cóncavos en la caja de masas pequeñas
(Hessiana diagonal de sitio más los términos positivos de línea para $-H$).
La continuación exterior y esta silla identifican el predictor reoptimizado de R3.
No se ha tomado un máximo con actividades complejas.

\paragraph{Normalización exacta.}
Para una fibra $f$, sea $C_f=\diag(x_f)-x_fx_f^t$ su covarianza categórica.
Se tiene $\det C_f=q_f\prod_{i\in f}x_i$. Al separar las masas marcadas, la acción
local es $H_B+R$ y Schur da, además de $A_*$ y la constante de componente,
el menos medio logaritmo de
\[
 D_{\rm loc}=\det\big[(-H_B''-R_c'') V_0\big]\prod_f
                     \frac{\det C_f}{\prod_{i\in f}v_i}.
\]
En efecto el producto escalar junto a $\det(-H_B''-R_c'')$ se reduce a
$\prod_i x_i(1-x_i)^{1-d_i}\prod_f q_f$. Los factores $q_f=t_b$ de RC
se compensan con $+\frac12\sum_b\log t_b$ en $A_*$; quedan los escalares
físicos del Schur de R3, 08b. La identidad es algebraica también para un ciclo.
La columna $i$ de $(-H_B''-R_c'')V_0$ se vuelve el vector unitario $e_i$ en $x_i=0$.
Se ve usando $-H_{\rm cat}''=\diag(1/x_f)+{\bf1}{\bf1}^t/q_f$ y la fórmula
de $H_B$. El bloque restante es el de la cara. Los cocientes escalares reducen
a $q_f/\prod_{i\in f}(1-x_i)$. El valor estacionario se escribe sin $\log x_i$
por la ecuación de campo. Por tanto el valor y el determinante son holomorfos en
los ejes y coinciden al desactivar las incidencias de una marca de actividad cero.

El exacto es el logaritmo de $\sum_X \chi_Y(X)\zeta^X e^{g(X)}$, sobre los átomos
compatibles en RC, con $\chi_Y$ las capacidades diagonales. En un radio fijo pequeño
la suma está cerca de uno. Los términos comunes a todos los $Y$ se quitan.
En el producto local sin reacción exterior se pone $R_c=A_c=g=0$.
Por lo anterior, el cambio del error desde este modelo local al físico es $O_v(s)$
en el polidisco: las sillas varían $O_v(s)$ y todos los determinantes regularizados
están cerca de uno. El prefactor exterior varía $O_v(s)$ por el apartado 2.
En la acción se puede sustituir $\log(\zeta_i/x_i)$ por su expresión holomorfa.
Cada diferencia proyectada se anula en $\zeta_i=0$ para todo $i\in S$. Si su
norma es $O_v(s^a)$ en el polidisco, dividir por $\prod_i\zeta_i$ y aplicar Cauchy
en un radio menor da $O_v(s^{a+v})$ en los campos físicos.

\paragraph{Bosques.}
Para un bosque de factores, las creencias categóricas se ensamblan exactamente
multiplicando por condicionadas desde las hojas. La misma construcción con
gaussianas de covarianzas $C_f$ muestra que el determinante local sin reacción
es uno. Con el campo desplazado por $\nabla R_c$, $V=(-H_B'')^{-1}$ es la
covarianza exacta del bosque. El término adicional de Schur es
$-\frac12\log\det(I-VR_c'')$. Con $m,V$ las medias y la covarianza a campo sin
desplazar, el cambio del predictor hasta primer orden en $s$ es
\[
 a[m]+\tfrac12 H[m,m]+\tfrac12\tr(VH).
\]
Por estación, el movimiento del optimizador de la parte independiente tiene
coste de segundo orden. El exacto da el mismo polinomio, por el apartado 2 y
$\E[XX^t]=V+mm^t$ en el bosque. Estos desarrollos son uniformes para las
actividades del polidisco. Por ello el error proyectado de un bosque cuyos
subcasos siguen siendo bosques es $O_v(j^{-v-2})$.
